% !TEX root =./ECE444_Practice_main.tex
%
% L3 practice set (Polarization and Bandwidth), ported from the Jupyter Book
% practice.md / practice-solutions.md into the exam class for Gradescope.
% Blank copy: \begin{solution}[h] reserves h of work space. SOLUTIONS copy
% (\iskey=1) prints the red worked solutions.

\fullwidth{\textbf{Learning Objective 1.3: } I can determine the polarization of
an antenna and describe the bandwidth characteristics of common antenna types.
\\[4pt]\normalfont\small Show your work. A \textbf{Documentation} line at the end
is required (write \textbf{None} if you did not collaborate).}

\begin{questions}

%%%%%%%%%%%%%%%%%%%%%%%%%%%%%%%%%%%%%%%%%
\question \textbf{Identify the polarization.} For each plane wave propagating in
$+\zhat$, identify the polarization (linear, RHCP, LHCP, or elliptical) and ---
where linear --- the tilt angle relative to $\xhat$.
\begin{parts}
	\part $\vec{E} = \xhat 3\cosp{\omega t - kz}$
		\begin{solution}[0.4in] Only an $\xhat$ component --- \textbf{linear along $\xhat$} ($0\mydeg$). \end{solution}
	\part $\vec{E} = \xhat 2\cosp{\omega t - kz} + \yhat 2\cosp{\omega t - kz}$
		\begin{solution}[0.4in] Equal amplitudes, in phase --- \textbf{linear, tilted $+45\mydeg$}. \end{solution}
	\part $\vec{E} = \xhat 5\cosp{\omega t - kz} + \yhat 5\cosp{\omega t - kz - 90\mydeg}$
		\begin{solution}[0.5in] Equal amplitudes, $\delta = -90\mydeg$, $+\zhat$ propagation --- \textbf{RHCP}. \end{solution}
	\part $\vec{E} = \xhat 5\cosp{\omega t - kz} + \yhat 5\cosp{\omega t - kz + 90\mydeg}$
		\begin{solution}[0.4in] $\delta = +90\mydeg$ --- \textbf{LHCP}. \end{solution}
	\part $\vec{E} = \xhat 3\cosp{\omega t - kz} + \yhat 1\cosp{\omega t - kz - 90\mydeg}$
		\begin{solution}[0.6in] Unequal amplitudes, $\delta = -90\mydeg$ --- \textbf{right-hand elliptical}, $\text{AR} = 3$ ($\approx 9.5\db$). \end{solution}
\end{parts}

\newpage
%%%%%%%%%%%%%%%%%%%%%%%%%%%%%%%%%%%%%%%%%
\question \textbf{Axial ratio.} A satellite downlink antenna is spec'd circularly
polarized with $\text{AR} \le 3\db$ across its operating band.
\begin{parts}
	\part Convert the $3\db$ axial ratio to a linear ratio (major/minor axis).
		\begin{solution}[0.9in] $\text{AR} = 10^{3/20} \approx 1.41$ (i.e.\ $\sqrt{2}$). \end{solution}
	\part Write the two orthogonal linear components (up to a common scale) with the
	appropriate $90\mydeg$ phase, assuming the ellipse major axis is along $\xhat$.
		\begin{solution}[1.1in]
		\eq{ E_x(t) = 1 \cdot \cosp{\omega t}, \qquad E_y(t) = 0.707\cosp{\omega t \pm 90\mydeg} }
		Sign of the phase sets the RHCP vs LHCP sense.
		\end{solution}
	\part An $\xhat$ receiver picks off $E_x$; a $\yhat$ receiver picks off $E_y$.
	Compute the ratio of received \emph{power} in dB. Should it match part (a)?
		\begin{solution}[1.2in]
		\eq{ \dfrac{P_x}{P_y} = \lp \dfrac{1.0}{0.707} \rp^{2} = 2.0 \longrightarrow 10\mylog{2.0} \approx 3.0\db }
		Yes --- the axial ratio in dB equals the received-power difference between the two aligned linear receivers.
		\end{solution}
\end{parts}

\newpage
%%%%%%%%%%%%%%%%%%%%%%%%%%%%%%%%%%%%%%%%%
\question \textbf{Polarization loss factor.}
\begin{parts}
	\part A ground station transmits a \emph{linearly} polarized signal along $\xhat$.
	A satellite receiver is \textbf{RHCP} and pointed correctly. Compute the PLF in dB.
		\begin{solution}[0.7in] Linear $\leftrightarrow$ CP: $\text{PLF} = 0.5 = -3\db$. \end{solution}
	\part Same ground station, but the satellite receiver is \textbf{LHCP}. Compute the PLF.
		\begin{solution}[0.7in] Still $-3\db$ --- linear-to-CP does not depend on sense. \end{solution}
	\part Two linear whip antennas: radio A vertical, radio B tilted $30\mydeg$ from
	vertical. Compute the PLF.
		\begin{solution}[0.8in] $\text{PLF} = \cos^{2}(30\mydeg) = (0.866)^{2} \approx 0.75 \longrightarrow 10\mylog{0.75} \approx -1.2\db$. \end{solution}
	\part In one sentence, explain why aircraft-to-ground links almost always specify
	circular polarization on at least one end.
		\begin{solution}[0.9in] Airframe attitude (roll/pitch/yaw) sweeps a linear aircraft antenna through every tilt relative to the ground station; making one end CP caps the worst-case polarization loss at $3\db$ regardless of attitude. \end{solution}
\end{parts}

\newpage
%%%%%%%%%%%%%%%%%%%%%%%%%%%%%%%%%%%%%%%%%
\question \textbf{Impedance bandwidth from a VSWR spec.} An antenna is spec'd
``VSWR $\le 2$ from $2.30\ghz$ to $2.55\ghz$.''
\begin{parts}
	\part What is the impedance bandwidth in MHz?
		\begin{solution}[0.5in] $2.55 - 2.30 = 250\mhz$. \end{solution}
	\part What is the fractional bandwidth (FBW) in percent?
		\begin{solution}[0.9in] \eq{ f_c = 2.425\ghz, \qquad \text{FBW} = \dfrac{250}{2425} \approx 10.3\% } \end{solution}
	\part Convert VSWR $= 2$ to $\left| \Gamma \right|$ and to return loss in dB.
		\begin{solution}[1in] \eq{ \left| \Gamma \right| = \dfrac{2-1}{2+1} = \dfrac{1}{3} \approx 0.333, \qquad \text{RL} = -20\mylog{0.333} \approx 9.5\db } \end{solution}
	\part Narrowband, broadband, or ultra-wideband? Likely family --- patch, horn, or
	log-periodic?
		\begin{solution}[0.9in] FBW $\approx 10\%$ is \textbf{broadband}: a standard-gain horn or a well-designed half-wave dipole. A patch would need matching tricks; a log-periodic is overkill. \end{solution}
\end{parts}

\newpage
%%%%%%%%%%%%%%%%%%%%%%%%%%%%%%%%%%%%%%%%%
\question \textbf{Fractional and ratio bandwidth.}
\begin{parts}
	\part A cellular antenna covers $698$--$960\mhz$ (low) and $1.71$--$2.69\ghz$
	(high). Compute the FBW and RBW of each band.
		\begin{solution}[1.7in]
		\eq{
			\text{Low: } & f_c = 829\mhz, \quad \text{FBW} = \dfrac{262}{829} \approx 31.6\%, \quad \text{RBW} = \dfrac{960}{698} \approx 1.38{:}1 \\
			\text{High: } & f_c = 2200\mhz, \quad \text{FBW} = \dfrac{980}{2200} \approx 44.5\%, \quad \text{RBW} = \dfrac{2690}{1710} \approx 1.57{:}1
		}
		\end{solution}
	\part A UWB radar antenna covers $3.1$--$10.6\ghz$. Compute its RBW. Does it meet
	the FCC UWB threshold of RBW $\ge 2{:}1$?
		\begin{solution}[0.9in] \eq{ \text{RBW} = \dfrac{10.6}{3.1} \approx 3.42{:}1 } Yes --- well above the $2{:}1$ threshold. \end{solution}
	\part Which antenna family would you pick for each of the three bands, with a
	one-sentence reason?
		\begin{solution}[1.4in] Low ($\approx 32\%$): broadband monopole or PIFA. High ($\approx 45\%$): chassis slot or stacked patch. UWB ($3.4{:}1$): Vivaldi (tapered slot) or log-periodic --- self-scaling over a wide band without matching tricks. \end{solution}
\end{parts}

\newpage
%%%%%%%%%%%%%%%%%%%%%%%%%%%%%%%%%%%%%%%%%
\question \textbf{Chu-Harrington intuition.} A design must fit inside a sphere of
radius $a = 15\text{ mm}$ and operate at $f = 915\mhz$ (ISM band).
\begin{parts}
	\part Compute $\lambda$, $k$, and the electrical size $ka$.
		\begin{solution}[1.2in] \eq{ \lambda = \dfrac{c}{f} \approx 0.328\m, \quad k = \dfrac{2\pi}{\lambda} \approx 19.2\ \text{rad/m}, \quad ka = 19.2(0.015) \approx 0.288 } \end{solution}
	\part Estimate the minimum $Q$ from $Q_{\min} \approx \dfrac{1}{(ka)^{3}} + \dfrac{1}{ka}$.
		\begin{solution}[1in] \eq{ Q_{\min} \approx \dfrac{1}{0.0239} + 3.47 \approx 45 } \end{solution}
	\part Estimate the maximum fractional bandwidth $\text{FBW}_{\max} \approx 1/Q_{\min}$
	(for VSWR $\le 2$).
		\begin{solution}[0.8in] $\approx 1/45 \approx 2.2\%$ ($\approx 20\mhz$ at 915 MHz). \end{solution}
	\part The designer wants 100 MHz of usable bandwidth at 915 MHz. Possible in a
	15 mm sphere?
		\begin{solution}[1.2in] No --- $100\mhz$ is FBW $\approx 11\%$, roughly $5\times$ the Chu-Harrington ceiling for a 15 mm sphere. It requires a larger antenna, or resistive loading / matching-network absorption that trades gain for bandwidth. \end{solution}
\end{parts}

\newpage
%%%%%%%%%%%%%%%%%%%%%%%%%%%%%%%%%%%%%%%%%
\question \textbf{Reading the datasheet.} A vendor offers a GPS patch specified as
\textbf{RHCP}, $\text{AR} \le 3\db$, ``bandwidth 4\% (VSWR $\le 2$)''. You are
choosing the antenna for the other end of the link.
\begin{parts}
	\part The datasheet defines handedness \emph{from the receiver looking back
	toward the incoming wave}. Which convention is that, what is the part's sense in
	\textbf{IEEE} terms, and what goes wrong if you pair it with a satellite
	transmitting IEEE RHCP?
		\begin{solution}[1.4in]
		That is the \textbf{optics/physics} convention, opposite to IEEE. The part is
		IEEE \textbf{LHCP}. Paired with an IEEE RHCP satellite it is a \emph{sense
		mismatch} --- ideally a null, and in practice only the few dB of rejection set
		by the real axial ratio (part (e)). This is a link-killer, not a few-dB penalty.
		\end{solution}
	\part Convert $\text{AR} = 3\db$ to a linear ratio $A$, and give $A^{2}$.
		\begin{solution}[0.8in] \eq{ A = 10^{3/20} \approx 1.41, \qquad A^{2} \approx 2.0 } \end{solution}
	\part A \emph{linear} antenna receives from this patch. Using
	$\text{PLF}(\theta) = \dfrac{A^{2}\cos^{2}\theta + \sin^{2}\theta}{A^{2}+1}$, find
	the best and worst PLF over all orientations $\theta$, in dB.
		\begin{solution}[1.5in]
		\eq{ p_{\max} = \dfrac{A^{2}}{A^{2}+1} \approx \dfrac{2}{3} \longrightarrow 10\mylog{0.667} \approx -1.8\db }
		\eq{ p_{\min} = \dfrac{1}{A^{2}+1} \approx \dfrac{1}{3} \longrightarrow 10\mylog{0.333} \approx -4.8\db }
		Not the flat $-3\db$ of the ideal table --- and you rarely control $\theta$.
		\end{solution}
	\part Show that the peak-to-null swing in dB equals the axial ratio in dB.
		\begin{solution}[1.3in]
		\eq{ \dfrac{p_{\max}}{p_{\min}} = \dfrac{A^{2}/(A^{2}+1)}{1/(A^{2}+1)} = A^{2} }
		\eq{ 10\mylog{A^{2}} = 20\mylog{A} = \text{AR}_{\text{dB}} }
		Here $10\mylog{2.0} \approx 3.0\db$, matching the spec. The result is exact, not
		a rule of thumb.
		\end{solution}
	\part The far end is an opposite-sense CP antenna, also $\text{AR} = 3\db$, major
	axes aligned. Worst-case rejection is $\lp \dfrac{A^{2}-1}{A^{2}+1} \rp^{2}$.
	Evaluate it and compare with the idealized PLF table.
		\begin{solution}[1.3in]
		\eq{ \lp \dfrac{2-1}{2+1} \rp^{2} = \dfrac{1}{9} \approx 0.111 \longrightarrow 10\mylog{0.111} \approx -9.5\db }
		(Carrying $A^{2} = 1.995$ instead of $2$ gives $-9.6\db$.) The ideal table says
		$-\infty\db$. Real hardware gives roughly $10\db$ of cross-pol isolation ---
		useful, but nowhere near a null.
		\end{solution}
	\part Name two things you must ask before believing the antenna works over the
	advertised 4\%.
		\begin{solution}[1.3in]
		\textbf{Which} bandwidth --- 4\% is the \emph{impedance} BW at VSWR $\le 2$,
		while the axial-ratio and pattern bandwidths are narrower and go unquoted ---
		and \textbf{over what scan angle}, since AR is normally specified at boresight
		only and degrades off-axis. A defensible third: at what VSWR threshold, since
		the number moves with the bar.
		\end{solution}
\end{parts}


\vspace{1.5em}
\noindent\textbf{Documentation:} \rule[-2pt]{0.6\linewidth}{0.4pt}

\end{questions}
